\chapter{Chapter Outline}\label{ch:chapter_outline}
\markboth{Chapter Outline}{}

% ==========================================================================
% This chapter is NOT part of the final thesis. It serves as a planning
% reference that summarizes what each chapter and section should contain.
% Remove the \input for this file from FAPS-Thesis.tex before submission.
% ==========================================================================

% ------------------------------------------------------------------
\section{Chapter 1: Introduction (~5 pages)}
% ------------------------------------------------------------------

\subsection{Motivation and Context}
\begin{itemize}
    \item Introduce the importance of autonomous robotic cloth manipulation in manufacturing and recycling contexts
    \item Highlight the challenges of handling deformable objects (textiles) and why simulation with reinforcement learning is a promising approach
    \item Emphasize the specific scenario: using a 6-DOF tendon-driven robotic arm to pick up and spread cloth for inspection on a conveyor
    \item Establish the industrial relevance: textile sorting, quality inspection, and recycling as concrete application domains
    \item Briefly frame the sim-to-real challenge as a central theme connecting simulation training to physical deployment
\end{itemize}

\subsection{Problem Statement and Objectives}
\begin{itemize}
    \item Clearly state the thesis problem: developing an RL-based control strategy for a tendon-driven robot to manipulate cloth (grasp, stretch, present) in simulation and eventually real life
    \item List the research objectives:
    \begin{enumerate}
        \item Establishing a high-fidelity Isaac Sim simulation environment
        \item Designing and training RL policies for cloth manipulation tasks
        \item Integrating a custom tendon-driven robot model
        \item Addressing sim-to-real transfer (ensuring the learned policy can generalize or be adapted)
    \end{enumerate}
    \item Explain that initial simpler manipulation tasks have already been implemented as a stepping stone and the thesis builds on this foundation
    \item Scope delimitation: explicitly state what is NOT addressed (e.g., full autonomous perception pipeline, multi-robot coordination) and justify why
\end{itemize}

\subsection{Thesis Outline}
\begin{itemize}
    \item Summarize the structure of the thesis (Chapters 2--7) for the reader
    \item Chapter 2: background theory and related work
    \item Chapter 3: research gap and refined problem
    \item Chapter 4: methodology including simulation setup and RL approach
    \item Chapter 5: experimental results
    \item Chapter 6: discussion of findings and limitations
    \item Chapter 7: conclusion and future work
\end{itemize}

% ------------------------------------------------------------------
\section{Chapter 2: Background and State of the Art (~20 pages)}
% ------------------------------------------------------------------

\subsection{Section 2.1: Reinforcement Learning for Robotic Manipulation}
\begin{itemize}
    \item Introduce Markov Decision Processes (MDP) as the formalism for RL: define states, actions, transitions, rewards, and discount factor
    \item Outline key RL algorithms relevant to continuous robot control: model-free policy gradient methods (DDPG, PPO, SAC)
    \item Discuss why these algorithms are suitable for high-dimensional, continuous action problems like controlling robotic arms
    \item Compare on-policy (PPO) vs.\ off-policy (SAC, DDPG) approaches: sample efficiency, stability, and suitability for GPU-parallel simulation
    \item Highlight specific algorithms or frameworks used in this thesis (e.g., RL Games, Stable Baselines, skrl)
    \item Briefly cover policy network architectures: MLPs, actor-critic structure, and optional recurrent layers for partial observability
    \item Summarize prior successes of RL in simulation for manipulation tasks (e.g., OpenAI Dexterous Hand, QT-Opt)
\end{itemize}

\subsection{Section 2.2: Physics Simulation in Isaac Sim 5.1}
\begin{itemize}
    \item Introduce the physics engine stack of Isaac Sim 5.1: NVIDIA PhysX 5 as the core solver, with GPU-accelerated simulation capabilities
    \item Explain why understanding the underlying physics models is essential for sim-to-real transfer and RL training fidelity
\end{itemize}

\subsubsection{Rigid Body Dynamics}
\begin{itemize}
    \item Cover how Isaac Sim handles rigid body simulation via PhysX: articulated bodies, joints, contacts, and friction models
    \item Explain the Temporal Gauss-Seidel (TGS) solver used by PhysX for rigid body constraint solving
    \item Describe relevant simulation parameters: time step, solver iterations, contact offsets, friction coefficients
    \item Discuss GPU-accelerated rigid body simulation for parallel RL environments
    \item Explain how articulated robots (including the tendon-driven arm and gripper) are simulated as rigid body chains with joint constraints
    \item Cover collision detection: mesh vs.\ convex decomposition vs.\ primitive shapes; performance vs.\ accuracy tradeoffs
\end{itemize}

\subsubsection{Cloth and Deformable Object Simulation}
\begin{itemize}
    \item Explain why deformable objects break standard rigid-body assumptions: applied forces not only translate the object but also change its shape and contact state
    \item Describe common cloth modeling approaches in context: mass-spring models, particle-based constraint methods, continuum formulations (FEM/MPM)
    \item Emphasize the unique challenges of cloth dynamics: high degrees of freedom, self-collision, sensitivity to material parameters (bending stiffness, friction, damping)
\end{itemize}

\subsubsection{Position-Based Dynamics (PBD) --- PhysX Particle Cloth}
\begin{itemize}
    \item Explain PBD as the solver used by PhysX for particle cloth: particles with mass connected via holonomic constraints
    \item Describe the PBD time-step algorithm: external-force velocity update, position prediction, iterative Gauss-Seidel constraint projection, velocity correction
    \item Present the constraint correction formula and explain how stiffness depends on time step and iteration count
    \item Introduce XPBD as an extension with compliance formulation and Lagrange multipliers for more physically meaningful stiffness
    \item Explain why PhysX uses PBD for particle cloth: GPU parallelism via graph coloring, robustness, interaction with rigid bodies and articulations
    \item Discuss how PBD particle cloth is characterized in Omniverse/Isaac Sim: a PBD mass-spring system distinct from newer deformable formulations
\end{itemize}

\subsubsection{Vertex Block Descent (VBD) --- Newton Thin Deformables}
\begin{itemize}
    \item Explain VBD as an implicit, energy-based solver targeting variational implicit Euler time integration
    \item Describe the global energy minimization formulation and the vertex-level block coordinate descent approach
    \item Present the local Newton update step ($3 \times 3$ Hessian solve per vertex)
    \item Discuss key advantages: unconditional stability, convergence to implicit Euler solution, GPU parallelizability via vertex coloring
    \item Explain Newton's role in NVIDIA's robotics context: improved performance and stability of VBD for thin deformables such as clothing
    \item Discuss why VBD is attractive for cloth manipulation: better tolerance of stiffness ratios (near-inextensible stretch with softer bending) and larger time steps
\end{itemize}

\subsubsection{Comparison: PBD vs.\ VBD for Cloth in RL Workflows}
\begin{itemize}
    \item Compare mathematical targets: PBD projects positions to satisfy constraints vs.\ VBD minimizes variational implicit Euler energy
    \item PBD: simple, robust, fast, but effective stiffness depends on time step and iteration count; harder to map to physical material properties
    \item VBD: principled implicit integration, strong stability guarantees, favorable convergence in large-DOF systems
    \item Relate to Omniverse ecosystem: particle cloth = PBD mass-spring; newer deformable features moving toward implicit/compliance-based formulations
    \item Motivate the choice for this thesis: PhysX/PBD particle cloth for throughput and robustness in Isaac Sim 5.1; Newton/VBD when greater stability under stiff cloth and larger time steps is needed
\end{itemize}

\subsection{Section 2.3: NVIDIA Isaac Sim and IsaacLab Architecture}
\begin{itemize}
    \item Introduce NVIDIA Isaac Sim (version 5.1) as the simulation environment
    \item Describe core features: built on Omniverse, high-fidelity physics (PhysX 5), photorealistic rendering with RTX, ROS 2 support, GPU acceleration for parallel simulations
    \item Explain IsaacLab: NVIDIA's unified robotics learning framework built on Isaac Sim (evolved from the Orbit framework)
    \item Note that IsaacLab combines Isaac Gym's GPU-accelerated physics with Orbit's environment library, providing a modular toolkit for RL, imitation, and motion planning
    \item Discuss how IsaacLab's MDP workflow organizes simulation tasks: defining robots, environments, and training pipelines in a reusable way
    \item Clarify why IsaacLab + Isaac Sim 5.1 is well-suited for this thesis (integrated cloth dynamics support, tools for distributed RL training)
    \item Compare with alternative simulation platforms (MuJoCo, PyBullet, SOFA) to justify the choice of Isaac Sim for this specific use case
\end{itemize}

\subsection{Section 2.4: Tendon-Driven Robot Kinematics and Modeling}
\begin{itemize}
    \item Explain what tendon-driven robots are: actuators apply force via tendons (cables) to move joints, allowing lighter, more compliant designs but requiring complex modeling
    \item Cover basic kinematic principles: how pulling tendons causes joint rotation or continuous bending, constant-curvature assumptions, differences in inverse kinematics/control from rigid joint robots
    \item Discuss the specific robot model relevant to the project: the 2-DOF linear stage + 3-DOF tendon arm + Robotiq gripper assembly
    \item Describe challenges in simulating tendon-driven systems in Isaac Sim: modeling tendon routing and elasticity, approximation via constraint forces or URDF/MJCF model with mimic joints
    \item Include literature context on tendon-driven manipulators: how their precision and dynamics differ from direct-drive robots
    \item Discuss the advantages of tendon-driven designs for compliant cloth interaction (inherent compliance, force limiting, safe contact)
    % Workspace analysis:
    \item Introduce the concept of workspace analysis for serial manipulators: the set of all end-effector poses reachable within joint limits
    \item Define forward kinematics (FK) as the mapping from joint configuration $q$ to end-effector position/orientation; explain why FK-based sampling is used when analytical workspace boundaries are intractable (high DOF, non-standard kinematics)
    \item Present the Monte Carlo sampling approach: uniformly sample joint configurations within joint limits, compute FK for each sample, and aggregate results in Cartesian voxels (analogous to MathWorks \texttt{generateRobotWorkspace})
    \item Define the \textbf{Yoshikawa manipulability index} $w(q) = \sqrt{\det(J(q)\,J(q)^\top)}$: measures how well the end effector can be moved in all Cartesian directions; zero at singularities, higher values indicate better dexterity [Yoshikawa 1985]
    \item Define the \textbf{inverse condition number} (isotropy index) $\kappa_{\mathrm{inv}} = \sigma_{\min}/\sigma_{\max}$ of the geometric Jacobian: ranges from 0 (singular) to 1 (perfectly isotropic); provides a complementary measure of motion quality [Salisbury \& Craig 1982]
    \item Define \textbf{reachability density}: the number of FK samples per voxel, proportional to the joint-space volume mapping to that Cartesian region; higher density indicates more configurations can reach that area
    \item Discuss how these measures inform task-space design: placing objects and setting action bounds in regions with high manipulability and reachability, avoiding singular zones
\end{itemize}

\subsection{Section 2.5: Policy Learning Pipeline (MDP Workflow)}
\begin{itemize}
    \item Detail the standard RL pipeline in the context of IsaacLab's MDP workflow
    \item Define MDP components for the cloth manipulation task:
    \begin{itemize}
        \item State (observations)
        \item Action space (robot joint torques/velocities or high-level gripper motions)
        \item Reward function design (what constitutes success for grasping/stretching cloth)
        \item Episode termination conditions
    \end{itemize}
    \item Explain curriculum learning or reward shaping concepts if used
    \item Describe how IsaacLab provides API/structure for tasks (environment wrappers, vectorized environments on GPU)
\end{itemize}

\subsection{Section 2.6: Prior Work in Robotic Cloth Manipulation}
\begin{itemize}
    \item Review existing research on robotic cloth/deformable object manipulation
    \item Cover both model-free approaches (deep RL, learning-based methods) and model-based/analytical approaches (vision-based controllers, heuristic methods)
    \item Cite key studies and benchmarks: folding/unfolding tasks, cloth smoothing/flattening (Seita et al.), GarmentBench/GarmentLab benchmarks
    \item Note outcomes: model-free RL achieves good simulation performance but needs many training samples and can struggle to generalize
    \item Highlight open problems from surveys: generalizing to varying fabrics, dealing with noisy perception, accomplishing long-horizon multi-step cloth tasks
    \item Discuss sim-to-real approaches specifically for deformable objects: domain randomization of cloth parameters, visual domain adaptation, teacher-student policy distillation
    \item Situate the thesis in context: how it builds on prior work and what gaps it addresses
    \item Provide a structured comparison table of key related works (method, robot type, task, simulation used, sim-to-real status)
\end{itemize}

\subsection{Section 2.7: Robot and Scene Description Formats (USD, URDF, MJCF)}
\begin{itemize}
    \item Explain USD (Universal Scene Description): extensible 3D scene graph format for constructing simulation scenes in Isaac Sim (objects, lighting, physics properties)
    \item Explain URDF (Unified Robot Description Format) and MJCF (MuJoCo XML) as common robot model formats
    \item Clarify which format is used for the tendon-driven robot and why
    \item Explain the role of each format: URDF/MJCF capture kinematic structure, joint limits, inertial properties; USD integrates robot and cloth into the simulated world
    \item Note any modifications to the robot's URDF to model tendons (workarounds for standard URDF limitations)
\end{itemize}

% ------------------------------------------------------------------
\section{Chapter 3: Problem Statement and Research Gap (~3 pages)}
% ------------------------------------------------------------------

\subsection{Section 3.1: Critical Analysis of State of the Art}
\begin{itemize}
    \item Build directly on Chapter 2's review and summarize what has been done and where limitations remain
    \item Note that while prior works achieve either cloth manipulation with traditional robots or basic RL policies in simulation, the combination of a tendon-driven arm with complex cloth tasks is unexplored
    \item Point out unresolved questions:
    \begin{itemize}
        \item How to effectively train an RL policy for a high-DOF tendon-driven manipulator on cloth tasks without extensive manual reward shaping?
        \item How to maintain sim-to-real fidelity for cloth interactions and tendon dynamics?
    \end{itemize}
\end{itemize}

\subsection{Section 3.2: Research Gap and Justification}
\begin{itemize}
    \item Clearly articulate the Handlungsbedarf (need for action) per FAPS guidelines
    \item Key gap: most existing studies use rigid-joint robots or simplified simulations; literature lacks demonstrations of RL-based control on tendon-driven robots for cloth handling; sim-to-real strategies for deformable tasks are not well established
    \item Justify scientific value: contributing knowledge on tendon-driven robot control and RL in deformable object tasks
    \item Justify practical relevance: automating textile recycling or clothing handling tasks
    \item Include a summary figure or diagram visualizing the research gap: map existing work along axes (robot type vs.\ object type vs.\ method) and show the empty quadrant this thesis fills
\end{itemize}

\subsection{Section 3.3: Refined Problem Definition}
\begin{itemize}
    \item Formulate precise research questions or hypotheses:
    \begin{itemize}
        \item How can a simulation environment be configured to realistically train cloth manipulation policies for a tendon-driven robot?
        \item Which RL algorithms and reward formulations yield stable learning for the cloth stretching task?
        \item What strategies can improve transferability of the learned policy to a physical robot (using ROS 2)?
    \end{itemize}
    \item Define the scope: specific cloth size/type, specific robot design, what is not addressed (e.g., vision-based perception may be simplified)
    \item List the main thesis contributions:
    \begin{enumerate}
        \item Integration of a tendon-driven robot model in Isaac Sim
        \item An RL pipeline for cloth grasp-and-stretch
        \item Evaluation of policy performance and generalization
    \end{enumerate}
\end{itemize}

\subsection{Section 3.4: Success Criteria (Optional)}
\begin{itemize}
    \item Define how success is measured: e.g., cloth spreading coverage in simulation, stability of tendon arm control, comparisons to baseline methods
    \item Tie the problem statement to the forthcoming methodology
\end{itemize}

% ------------------------------------------------------------------
\section{Chapter 4: Methodology and Implementation (~15--18 pages)}
% ------------------------------------------------------------------

\subsection{Section 4.1: Overview of Approach}
\begin{itemize}
    \item Provide a high-level overview of the solution strategy
    \item Include a diagram of the system architecture or workflow (IsaacLab environment loop: state $\rightarrow$ policy $\rightarrow$ action $\rightarrow$ environment)
    \item Describe development stages:
    \begin{enumerate}
        \item Stage 1: set up a simple simulation with a generic robot and train basic RL tasks
        \item Stage 2: integrate the tendon-driven robot model
        \item Stage 3: tackle the cloth manipulation task with the full system
    \end{enumerate}
\end{itemize}

\subsection{Section 4.2: Simulation Environment Setup}
\begin{itemize}
    \item Detail the Isaac Sim 5.1 environment configuration for training
    \item Describe the USD scene: table/conveyor, cloth object (size, mesh/primitive, physical properties like mass, friction, stiffness), initial conditions (crumpled/folded/randomized)
    \item Robot placement and any sensors in simulation
    \item If ROS 2 is used within Isaac Sim (ROS bridge extension), explain how ROS topics are configured
    \item Simulation parameters: physics time-step, faster-than-real-time execution for RL training
    \item Domain randomization strategy: which parameters are randomized (cloth mass, friction, stiffness, initial pose, lighting) and their ranges
    \item Include a figure of the complete simulation scene with labeled components
\end{itemize}

\subsection{Section 4.3: Robot Model Integration (Tendon-Driven Arm)}
\begin{itemize}
    \item Explain how the 6-DOF tendon-driven robotic arm is modeled and imported
    \item Describe the 2-DOF articulated base (linear Y--Z stage)
    \item Describe the 3-DOF tendon-driven arm
    \item Describe the Robotiq 2F-140 gripper integration
    \item Explain the modular assembly via Isaac Sim's Robot Assembler
    \item Describe any verification of the robot model in simulation (open-loop tests, joint motion verification)
    \item Calibration of tendon tensions, spring constants, joint limits, and stiffness
    \item Control gain tuning and validation
    \item Include a figure of the robot model
\end{itemize}

\subsubsection{Robot Simulation Setup and Validation}
\begin{itemize}
    \item Detail the physics simulation configuration for the articulated robot: articulation solver type (TGS), solver iteration counts, physics time step, and their effect on joint stability
    \item Describe the tendon actuation modelling decision: why PhysX articulation tendons were not used (fidelity and stability concerns per NVIDIA docs) and how direct PD joint drives approximate tendon behaviour instead
    \item Explain the PD drive parameterisation: stiffness and damping per joint, how initial values were derived from estimated tendon elasticity, and the iterative tuning process using Isaac Sim's Gain Tuner
    \item Describe collision geometry choices: convex hull generation from CAD meshes, performance vs.\ accuracy trade-off, and verification that no self-interpenetration occurs across the full joint range
    \item Present the step-response validation methodology: sinusoidal and step position commands applied to each joint independently; metrics recorded (settling time, overshoot, steady-state error); comparison with Klein's (2023) PID controller parameters and Jacobian transpose mapping
    \item Explain gravity compensation handling for the ceiling-mounted configuration and its effect on the vertical axis tuning
    \item Include a figure showing the Gain Tuner output: commanded vs.\ actual joint trajectories before and after tuning
\end{itemize}

\subsubsection{Workspace Analysis Methodology}
\begin{itemize}
    \item Describe the Monte Carlo FK sampling procedure: 200\,000 uniform random joint configurations within joint limits, FK computed via the IsaacSim 5.1 physics engine (not analytical FK, to capture the exact simulated kinematics)
    \item Explain parallelisation: 4096 environments running simultaneously on GPU, each evaluating independent joint configurations per step
    \item Detail the data collected per sample: end-effector position (Cartesian), geometric Jacobian, Yoshikawa manipulability index, inverse condition number
    \item Describe voxelisation of the Cartesian workspace for density and quality-metric aggregation
    \item Explain collision handling: scene objects (conveyors, drums) disabled during sampling so the kinematic envelope is fully explored; robot self-collision retained
    \item Define the desired workspace region derived from Klein's (2023) task geometry and how it is compared against the sampled reachable workspace
    \item Describe the visualisation pipeline: 3D scatter plots, 2D cross-section heatmaps through the workspace midpoint, colour-coded by manipulability and reachability density
    \item Explain how the results inform task-space bounds and action limits for the RL environment configuration
\end{itemize}

\subsection{Section 4.4: Reinforcement Learning Setup}
\begin{itemize}
    \item \textbf{State Space:} what observations the agent receives (joint angles/velocities, gripper status, cloth features/key points, camera images if used)
    \item \textbf{Action Space:} what actions the agent outputs (continuous joint commands: position increments, velocity, or torque); clarify how tendon coupling affects the action space
    \item \textbf{Reward Function:} reward design for cloth manipulation (e.g., negative distance from desired flat configuration, increase in covered area); intermediate shaping rewards (grasping, lifting); reward schedule or curriculum
    \item \textbf{RL Algorithm and Training Parameters:} which algorithm is used (e.g., PPO); library (RL Games or Stable Baselines); key hyperparameters (learning rate, discount factor, policy network architecture, training duration); number of parallel simulation instances for GPU-accelerated training; any pre-training or demonstrations
    \item \textbf{MDP Workflow in IsaacLab:} how the custom task class is defined, specifying observation and reward functions; IsaacLab wrappers for interfacing with the RL algorithm; role of ROS 2 during training (if any)
    \item \textbf{Reproducibility:} provide a complete hyperparameter table and specify exact software versions (Isaac Sim, IsaacLab, CUDA, Python, RL library) for full reproducibility per FAPS methodology requirements
\end{itemize}

\subsection{Section 4.5: Initial Experiments (Basic Task Policies)}
\begin{itemize}
    \item Describe initial RL experiments with simpler tasks and robots to validate the RL pipeline
    \item \textbf{Reach Pose Task:} setup, purpose, and what was learned/adjusted
    \item \textbf{Simple Pick and Place Task:} a simpler robot model (standard manipulator), dummy object instead of flexible cloth
    \item Purpose: test the training setup, reward shaping, hyperparameters before introducing tendon robot and deformable cloth
    \item Note adjustments from these initial trials (missing observations, hyperparameter tuning)
\end{itemize}

\subsection{Section 4.6: Advanced Task Implementation (Cloth Grasping and Stretching)}
\begin{itemize}
    \item Detail the final task scenario: tendon-driven robot grasps cloth from surface and stretches it out
    \item Break down the complex task into phases: grasp, lift, stretch, place
    \item Strategies for making learning feasible:
    \begin{itemize}
        \item Curriculum learning (start with cloth already grasped, or progressively increase difficulty)
        \item Auxiliary rewards for sub-goals (reward for grasping, reward for moving upward)
    \end{itemize}
    \item Gripper simulation: how grasp detection is handled (contact sensors, constraint on gripper close)
    \item Safety/stability measures: limiting force on tendons to realistic ranges to prevent simulation instabilities
\end{itemize}

\subsection{Section 4.7: Software and Hardware Integration (ROS 2)}
\begin{itemize}
    \item Describe how ROS 2 is incorporated, especially with an eye towards real-world deployment
    \item During simulation: ROS 2 as bridge (echoing states, visualization in RViz)
    \item Outline how the trained policy interfaces with a real robot via ROS 2:
    \begin{itemize}
        \item Robot's joint controllers subscribe to commanded positions from the RL policy
        \item Policy subscribes to sensor feedback (joint states, possibly depth camera for cloth observation)
    \end{itemize}
    \item Any preliminary setup for sim-to-real: using ROS 2 control for the tendon robot in sim to mirror what would be done on hardware
    \item Exporting the trained policy for a ROS 2 node
\end{itemize}

% ------------------------------------------------------------------
\section{Chapter 5: Results (~15 pages)}
% ------------------------------------------------------------------

\subsection{Section 5.1: Preliminary Experiments (Baseline Results)}
\begin{itemize}
    \item Report performance of the RL policy on simple pick-and-place or reaching tasks with the simplified robot
    \item Include learning curves: reward over training iterations for the baseline task
    \item Number of episodes needed to reach a success criterion
    \item If multiple algorithms or settings were tried, compare briefly (e.g., PPO vs.\ DDPG in reaching task)
    \item Purpose: demonstrate that the RL pipeline works on a simpler problem and record insights
    \item Justify choices for the harder task (confirm hyperparameters or algorithm selection)
\end{itemize}

\subsection{Section 5.2: Performance of Tendon-Driven Robot on Simple Tasks}
\begin{itemize}
    \item Present results if the tendon-driven arm was tested on a simple task to validate its controllability
    \item Report differences due to the robot's dynamics: whether learning was slower or required dealing with more complex dynamics
    \item Quantitative metrics: success rate, final accuracy for basic manipulation
    \item Validation of the simulated model: joint motions correspond to expected behavior, open-loop tests confirm model accuracy
\end{itemize}

\subsubsection{Robot Simulation Validation Results}
\begin{itemize}
    \item Present step-response plots for each joint: commanded vs.\ actual position over time, showing settling time, overshoot, and steady-state tracking
    \item Report quantitative tuning outcomes per joint: final stiffness and damping values, settling time ($\sim 0.1$--$0.3$\,s), overshoot ($< 5\%$), and steady-state error
    \item Compare the tuned PD behaviour with Klein's (2023) PID controller results: discuss differences due to the simplified tendon model and their practical implications
    \item Show that the vertical axis (base Z) required higher damping to counteract gravity-induced oscillations; wrist joints use lower stiffness reflecting tendon compliance
    \item Validate collision geometry: confirm no interpenetration across the full joint range, including extreme configurations
    \item Discuss any simulation artefacts observed during tuning (e.g., jitter at joint limits, contact instabilities at high stiffness) and how they were resolved
\end{itemize}

\subsubsection{Workspace Analysis Results}
\begin{itemize}
    \item Present the reachable workspace envelope: 3D scatter plot and volume estimate; compare against the desired workspace region from Klein's (2023) task geometry
    \item Show 2D cross-section heatmaps (e.g., Y--Z and X--Z slices through workspace midpoint) colour-coded by reachability density and manipulability
    \item Report Yoshikawa manipulability index distribution: mean, percentiles, and spatial variation; identify high-dexterity and near-singular regions
    \item Report inverse condition number distribution: highlight regions with isotropic vs.\ anisotropic motion capability
    \item Discuss coverage: what fraction of the desired task workspace is reachable, and at what quality level (e.g., ``95\% of the target volume is reachable with $w > 0.01$'')
    \item Identify practical implications: regions where the robot operates well vs.\ zones to avoid in task design; how workspace limits constrain object placement and RL action bounds
    \item Compare briefly with the workspace reported by Klein (2023) for the original arm design (without the 2-DOF base) to show the benefit of the linear stage
\end{itemize}

\subsection{Section 5.3: Cloth Manipulation Task Results}
\begin{itemize}
    \item \textbf{Learning Curve:} plot of reward vs.\ training steps or success rate vs.\ episodes; discuss how quickly the agent learned and whether it plateaued
    \item \textbf{Quantitative Performance:} define success metric (e.g., percentage of cloth spread flat above threshold); report final policy performance averaged over many test episodes (e.g., ``85\% cloth coverage on average, with $\geq$90\% coverage in 70\% of trials''); use tables and graphs
    \item \textbf{Policy Behaviors:} describe what the learned policy actually does (e.g., grasps corner, lifts, shakes, lays down stretched); illustrate with sequential images from simulation showing keyframes
    \item \textbf{Ablation or Variant Results:} different reward formulations, disabling certain observations; e.g., ``without tactile feedback, success rate dropped by 15\%''
    \item \textbf{Computational Cost:} report wall-clock training time, GPU memory usage, number of environment steps; quantify speedup from GPU parallelism
    \item \textbf{Sim-to-Real Considerations:} tests with added noise to gauge robustness; effect of domain randomization on policy performance; preliminary tests on real robot if available
\end{itemize}

\subsection{Section 5.4: Summary of Key Results (Optional)}
\begin{itemize}
    \item Concise bullet-point recap of the most important findings
    \item Answer the objectives as evidenced by results:
    \begin{itemize}
        \item Achieved reliable cloth picking in simulation
        \item Partial/full stretching accomplished (flatness percentage)
        \item RL training converged in approximately N million steps
    \end{itemize}
\end{itemize}

% ------------------------------------------------------------------
\section{Chapter 6: Discussion (~5 pages)}
% ------------------------------------------------------------------

\subsection{Section 6.1: Evaluation of Outcomes}
\begin{itemize}
    \item Discuss to what extent the thesis objectives (from Introduction) were achieved
    \item For each objective: was a robust policy learned? How effective is it?
    \item Provide reasoning for the performance observed (e.g., learning curve plateau indicates suboptimal but acceptable strategy)
    \item Connect with theory: high-dimensional cloth state making dense reward difficult, exploration challenges
\end{itemize}

\subsection{Section 6.2: Comparison with Related Work}
\begin{itemize}
    \item Compare approaches and results with literature from Section 2.6
    \item Discuss how the approach compared: advantages of GPU simulation enabling more training, using tendon-driven robot vs.\ standard robots
    \item Analyze shortcomings: e.g., pure RL approach vs.\ imitation learning may require additional strategies for real-world deployment
    \item Situate the contribution: either showing improvement or highlighting differences
\end{itemize}

\subsection{Section 6.3: Limitations of the Study}
\begin{itemize}
    \item Limitations in simulation fidelity: cloth simulation as approximation (PhysX may not capture all fine wrinkles), tendon model may not include cable elasticity perfectly
    \item Risk of policy overfitting to simulated environment quirks
    \item Generalization concerns: does the policy only work for a specific cloth size or initial pose? What about varying material properties?
    \item If ROS 2/hardware integration was not fully realized, state that results are limited to simulation
    \item Training time considerations and policy sensitivity to parameter changes
    \item Discuss the validity of chosen evaluation metrics: are they sufficient to judge real-world deployment readiness?
\end{itemize}

\subsection{Section 6.4: Practical Implications}
\begin{itemize}
    \item What the results mean for real-world applications (e.g., automated cloth handling in recycling facilities)
    \item Discuss the gap to reality: sensor noise, unpredictable cloth variations
    \item Whether the tendon-driven robot in simulation behaved similarly to expectations for reality; what adjustments might be needed
    \item Reflect on the decision to use IsaacLab/Isaac Sim: was it effective? Did GPU acceleration meaningfully speed up training?
\end{itemize}

\subsection{Section 6.5: Recommendations and Lessons Learned}
\begin{itemize}
    \item Key lessons from the development process, for example:
    \begin{itemize}
        \item Careful reward shaping was essential; purely sparse reward led to no learning
        \item Introducing incremental rewards at different task stages was crucial
        \item Calibration of joint limits and stiffness in simulation greatly affects training stability
        \item An improperly tuned model caused unstable training
    \end{itemize}
    \item These reflective points demonstrate critical thinking and provide valuable takeaways
\end{itemize}

% ------------------------------------------------------------------
\section{Chapter 7: Conclusion and Outlook (~3 pages)}
% ------------------------------------------------------------------

\subsection{Section 7.1: Conclusion (Summary of Findings)}
\begin{itemize}
    \item Recap the core accomplishments in a few concise paragraphs
    \item Reiterate the problem addressed and how it was solved
    \item Example summary: ``developed an RL framework for a tendon-driven robotic arm to manipulate cloth in simulation; built a high-fidelity Isaac Sim environment with cloth physics; integrated a custom tendon-driven robot model; using policy-gradient RL, the robot learned to pick up and stretch cloth''
    \item Highlight the significance: possibly one of the first demonstrations of RL on a tendon-driven robot for cloth manipulation, or a novel combination of IsaacLab with a custom robot model
    \item Touch on each objective from the introduction, confirming whether it was met
\end{itemize}

\subsection{Section 7.2: Outlook (Future Work)}
\begin{itemize}
    \item Specific next steps:
    \begin{enumerate}
        \item \textbf{Real-world implementation:} deploying the learned policy on the physical tendon-driven robot via ROS 2; evaluating on real cloth; dealing with sim-to-real gaps (domain randomization, fine-tuning with real data)
        \item \textbf{Improving the policy/generalization:} exploring more advanced RL algorithms; incorporating vision-based observation for handling varying cloth shapes or unknown cloth pieces
        \item \textbf{Task expansions:} moving beyond stretching to folding, handling multiple pieces, integrating higher-level planning for sequential tasks
        \item \textbf{Simulation enhancements:} more accurate cloth models, additional sensors (depth camera for cloth state perception)
    \end{enumerate}
    \item Open research questions: adaptability of learned policy to different fabrics (varying stiffness); training with a range of material properties for robustness
    \item Closing remark: reflect on broader impact --- integration of advanced simulation with RL demonstrates a path toward robots that can handle complex, deformable objects autonomously
\end{itemize}

% ------------------------------------------------------------------
\section{Appendices (not counted in page total)}
% ------------------------------------------------------------------

\begin{itemize}
    \item Supplementary material: extensive tables, hyperparameter lists, additional diagrams, code snippets
    \item Full URDF/MJCF of the robot
    \item Detailed network architectures
    \item Additional results plots
    \item Complete reward function formulations and their evolution during development
    \item Domain randomization parameter ranges table
    \item Software version matrix (Isaac Sim, IsaacLab, PhysX, CUDA, Python packages)
\end{itemize}
